\begin{definition}[Differential-form zero-degree scalar boundary]
\label{def:diffform-zero-degree-scalar-boundary}
Fix $\NameCert_{\ManifoldUp}$ and $\NameCert_{\TensorProductUp}$ sources
$M,T$ as in \autoref{def:diffform-bhist-form-carrier}. A differential-form
zero-degree scalar boundary row is a carried $\DiffFormUp$ row
$$
\omega_0=(d_0,\mathcal{P}_0,\tau_0,\sigma,\alpha_0,\lambda_0)
$$
accepted by $\mathsf{DiffFormBHistCarrier}_{M,T}$ with the following visible
coordinates:
\begin{enumerate}
  \item $d_0$ is the empty unary-zero history, so $\hsame(d_0,\emp)$ and
  $\mathsf{UnaryHistory}(d_0)$.
  \item $\mathcal{P}_0$ is the empty probe spine
  $\mathsf{ProbeBundle.Bnil}$.
  \item $\tau_0$ is the empty tensor-unit surface of the
  $\TensorProductUp$ source over $\mathcal{P}_0$.
  \item $\alpha_0$ is the empty antisymmetry ledger, with no adjacent-swap
  row and no cross-permutation row.
  \item $\sigma$ is the retained scalar classifier endpoint row supplied by
  the upstream manifold and tensor sources.
  \item $\lambda_0$ records exactly $M,T,d_0,\mathcal{P}_0,\tau_0,\sigma$,
  and $\alpha_0$ as the source ledger for this boundary row.
\end{enumerate}
Thus the row is a scalar-valued $0$-form surface inside the displayed
$\DiffFormUp$ carrier, not a host-level exception outside the carrier and
classifier of \autoref{def:diffform-bhist-form-carrier} and
\autoref{def:diffform-bhist-form-classifier}.
\end{definition}

\begin{theorem}[Differential-form zero-degree empty-probe exhaustion]
\label{thm:diffform-zero-degree-empty-probe-exhaustion}
Let $\omega_0$ be a differential-form zero-degree scalar boundary row of
\autoref{def:diffform-zero-degree-scalar-boundary}. For every probe row $p$,
the membership claim $\mathsf{InBundle}(p,\mathcal{P}_0)$ is impossible.
Equivalently, every downstream use of $\omega_0$ sees an empty probe spine,
with no hidden tangent input slot.
\end{theorem}
\begin{proof}
Unpack the boundary row through
\autoref{def:diffform-zero-degree-scalar-boundary}. Its probe coordinate is
the closed constructor $\mathsf{ProbeBundle.Bnil}$.

Invert an assumed $\mathsf{InBundle}(p,\mathcal{P}_0)$ claim. The
$\mathsf{ProbeBundle}$ membership constructors can only expose a head probe
from a nonempty spine or recurse into a nonempty tail. The constructor
$\mathsf{ProbeBundle.Bnil}$ has neither, so no membership constructor can
produce the assumed claim.

The same empty-spine inversion is compatible with the carrier and classifier
rows of \autoref{def:diffform-bhist-form-carrier} and
\autoref{def:diffform-bhist-form-classifier}. The scalar endpoint remains
present as $\sigma$, but no probe row can be read from the boundary.
\end{proof}

\begin{theorem}[Differential-form zero-degree scalar classifier readback]
\label{thm:diffform-zero-degree-scalar-classifier-readback}
Let
$$
\omega_0=(d_0,\mathcal{P}_0,\tau_0,\sigma,\alpha_0,\lambda_0)
$$
and
$$
\omega'_0=(d'_0,\mathcal{P}'_0,\tau'_0,\sigma',\alpha'_0,\lambda'_0)
$$
be two differential-form zero-degree scalar boundary rows over the same
$M,T$ sources. Then $\omega_0\sim_{\DiffFormUp}\omega'_0$ is exactly the
retained scalar endpoint classifier comparison $\sigma\sim_R\sigma'$,
together with the visible source-ledger comparisons
$$
\hsame(d_0,d'_0),\quad
\hsame(\tau_0,\tau'_0),\quad
\hsame(\alpha_0,\alpha'_0),\quad
\hsame(\lambda_0,\lambda'_0).
$$
No probe-indexed classifier clause remains after empty
$\mathsf{ProbeBundle.Bnil}$ inversion.
\end{theorem}
\begin{proof}
Read the two boundary rows through
\autoref{def:diffform-zero-degree-scalar-boundary}. Both probe spines are
$\mathsf{ProbeBundle.Bnil}$, both degree rows are the empty unary-zero
history, both tensor rows are empty tensor-unit surfaces, and both
antisymmetry ledgers are empty.

Now unfold the differential-form classifier of
\autoref{def:diffform-bhist-form-classifier}. The degree, tensor,
antisymmetry, and source-ledger clauses are exactly the displayed $\hsame$
comparisons on the boundary rows. The scalar endpoint clause is the retained
source classifier comparison $\sigma\sim_R\sigma'$.

It remains to rule out any further probe-indexed clause. By
\autoref{thm:diffform-zero-degree-empty-probe-exhaustion}, no
$\mathsf{InBundle}$ claim exists for either boundary spine. Hence every
matched multilinear probe row in the classifier is exhausted by the empty
tensor-unit endpoint, and the public readback surface is precisely the
scalar endpoint classifier plus the visible $\hsame$ ledger comparisons.
\end{proof}

\begin{theorem}[Differential-form zero-degree wedge Cont unit boundary]
\label{thm:diffform-zero-degree-wedge-cont-unit-boundary}
Let $\omega_0$ be a differential-form zero-degree scalar boundary row and let
$$
\eta=(d_\eta,\mathcal{P}_\eta,\tau_\eta,\rho,\beta,\mu)
$$
be any carried $\DiffFormUp$ row over the same $M,T$ sources. Whenever the
wedge rows $\omega_0\wedge_{\DiffFormUp}\eta$ and
$\eta\wedge_{\DiffFormUp}\omega_0$ are admitted by
\autoref{thm:diffform-wedge-ledger-coverage}, their probe and degree surfaces
are the nonempty side's surfaces:
$$
\mathsf{bundleAppend}(\mathsf{ProbeBundle.Bnil},\mathcal{P}_\eta)
\quad\text{and}\quad
\mathsf{bundleAppend}(\mathcal{P}_\eta,\mathsf{ProbeBundle.Bnil})
$$
are classified with $\mathcal{P}_\eta$, and the degree ledgers of
\autoref{def:diffform-wedge-degree-ledger} reduce through the $\Cont$ left
and right unit rows to $d_\eta$. The tensor-unit component of $\omega_0$
therefore contributes no additional probe slot and no additional degree
history to either wedge boundary.
\end{theorem}
\begin{proof}
First expose the wedge probe ledger. Apply
\autoref{def:diffform-wedge-probe-concatenation-ledger} and
\autoref{thm:diffform-wedge-ledger-coverage}. The output probe spine of the
left wedge is
$\mathsf{bundleAppend}(\mathsf{ProbeBundle.Bnil},\mathcal{P}_\eta)$, and the
output probe spine of the right wedge is
$\mathsf{bundleAppend}(\mathcal{P}_\eta,\mathsf{ProbeBundle.Bnil})$.

For the left wedge, invert the empty left spine. The
$\mathsf{bundleAppend}$ empty-left law identifies the output spine with
$\mathcal{P}_\eta$, and
\autoref{thm:diffform-zero-degree-empty-probe-exhaustion} rules out an
additional probe contributed by $\omega_0$.

For the right wedge, induct on $\mathcal{P}_\eta$. The empty case is the
same $\mathsf{ProbeBundle.Bnil}$ endpoint. In the cons case, the head probe is
preserved and the induction hypothesis identifies the appended tail with the
original tail, so the empty-right law identifies the output spine with
$\mathcal{P}_\eta$.

It remains to read the degree surface. By
\autoref{def:diffform-wedge-degree-ledger}, both admitted wedges carry
visible continuation rows for their output degrees. Since the boundary degree
of $\omega_0$ is $\emp$, the left wedge uses the $\Cont$ left-unit row
$\Cont(\emp,d_\eta,d_\eta)$, and the right wedge uses the $\Cont$ right-unit
row $\Cont(d_\eta,\emp,d_\eta)$. Thus the wedge output degree is classified
with $d_\eta$ on both sides.

Finally, the tensor component of $\omega_0$ is the empty tensor-unit surface
from \autoref{def:diffform-zero-degree-scalar-boundary}. The tensor
concatenation ledger named by
\autoref{thm:diffform-wedge-ledger-coverage} therefore records only the
unit-row concatenation with $\tau_\eta$, while the scalar endpoint comparison
remains in the retained source classifier described by
\autoref{def:diffform-bhist-form-classifier}. Repacking these probe, degree,
tensor, and endpoint rows gives the claimed unit boundary for both wedge
orders.
\end{proof}
